\clearpage
\hypersetup{pageanchor=false}
\setcounter{page}{1}
\maketitlesupplementary


\section{Extended Related Work}
\label{sec:suppl_related_work}

This section provides a comprehensive literature review of image-to-image translation, diffusion bridge models, and cross-modality translation to complement the specialized discussion on remote sensing in the main paper.

\paragraph{General Image-to-Image Translation.}
Image-to-image translation has been extensively studied in computer vision. Foundational works include Pix2Pix for paired conditional translation~\cite{isolaImageToImageTranslationConditional2017} and CycleGAN for unpaired translation~\cite{zhuUnpairedImageToImageTranslation2017}. Subsequent advances include contrastive learning for unpaired translation~\cite{parkContrastiveLearningUnpaired2020}, zero-shot~\cite{parmarZeroshotImagetoImageTranslation2023} and one-step translation with text-to-image models~\cite{parmarOneStepImageTranslation2024}.

\paragraph{Diffusion Bridge Models.}
Unlike standard diffusion models that generate from noise, diffusion bridge models learn a stochastic process that directly connects source and target distributions, making them well-suited for paired image-to-image translation. Denoising Diffusion Bridge Models (DDBM)~\cite{zhouDenoisingDiffusionBridge2024} define a variance-preserving bridge between data distributions. Related formulations include Brownian Bridge Diffusion Models (BBDM)~\cite{Li_2023_CVPR} and its bidirectional extension BiBBDM~\cite{xueBiBBDMBidirectionalImage2025}, Image-to-Image Schr\"odinger Bridge (I$^2$SB)~\cite{liuI2SBImagetoImageSchrodinger2023}, Diffusion Schr\"odinger Bridge matching~\cite{shiDiffusionSchrodingerBridge2023}, Dual Diffusion Implicit Bridges~\cite{suDualDiffusionImplicit2023a}, and Direct Diffusion Bridge with data consistency~\cite{chungDirectDiffusionBridge2023}. Recent extensions include Consistency Diffusion Bridge~\cite{heConsistencyDiffusionBridge2024}, Diffusion Bridge Implicit Models~\cite{zhengDiffusionBridgeImplicit2025}, deterministic translation via denoising Brownian bridge with dual approximators~\cite{xiaoDeterministicImagetoImageTranslation2025a}, and adaptive domain shift for cross-modality translation~\cite{wangAdaptiveDomainShift2026a}.

\paragraph{Cross-Modality Translation.}
Image translation across diverse sensor modalities---grayscale, near-infrared (NIR), thermal colorization, and color transfer---has been addressed with probabilistic color correction~\cite{oliveiraProbabilisticApproachColor2015}, NIR imagery colorization~\cite{suarezInfraRedImageryColorization2018}, deep colorization~\cite{zhangColorfulImageColorization2016,majumder2020deep}. 

\paragraph{GAN-based RS Translation.}
Early attempts at translating between SAR and optical modalities leveraged Generative Adversarial Networks (GANs). For instance, Wang et al.~\cite{wangSARtoOpticalImageTranslation2022} proposed a hierarchical latent feature approach to improve the fidelity of SAR-to-optical translation. Hu et al.~\cite{huGanbasedSAROptical2022} specifically addressed this task in fire-disturbed regions using GANs to reconstruct optical details from SAR signals, which is vital for post-fire damage assessment when clouds or smoke obscure optical sensors.

\paragraph{Diffusion Models for RS Translation.}
Reflecting the recent shift in generative modeling, diffusion models have been increasingly applied to RS translation tasks. Qin et al.~\cite{qinConditionalDiffusionModel2024} introduced a conditional diffusion model incorporating spatial-frequency refinement to capture both global structure and fine-grained textures in SAR-to-optical translation. To address the computational overhead of iterative sampling, Qin et al.~\cite{qinEfficientEndtoEndDiffusion2025} later developed an efficient end-to-end diffusion model capable of one-step SAR-to-optical translation, significantly reducing inference latency while maintaining high generation quality. Furthermore, Wang et al.~\cite{wangAdaptiveDomainShift2026a} proposed adaptive domain shift mechanisms within diffusion frameworks to better handle the large distribution gaps inherent in cross-modality remote sensing datasets. These advancements underscore the growing role of stochastic bridge processes and diffusion-based refinement in enabling robust, high-resolution translation across heterogeneous sensor modalities.


\section{EarthBridge Architecture and Implementation Details}
\label{sec:suppl_architecture}

This section provides the comprehensive technical details of the EarthBridge architecture, training hyperparameters, and multi-modal refinement modules.

\subsection{UNet Denoising Backbone}
Our denoiser $D_\theta$ is implemented in PyTorch~\cite{paszkePyTorchImperativeStyle2019} and based on a \texttt{UNet2DModel} backbone from the Hugging Face Diffusers library~\cite{vonplatenDiffusersStateoftheartDiffusion2022} with Adaptive Group Normalization (ADM-style~\cite{dhariwalDiffusionModelsBeat2021}). The detailed architectural parameters are summarized in \cref{tab:suppl_unet_arch}.

\begin{table}[h]
\centering
\caption{{\color{gray}UNet denoising backbone architectural parameters.}}
\label{tab:suppl_unet_arch}
{\color{gray}
\resizebox{\columnwidth}{!}{%
\begin{tabular}{ll}
\toprule
Property & Details \\
\midrule
Backbone Type & UNet2DModel (ADM-style~\cite{dhariwalDiffusionModelsBeat2021}) \\
Normalization & Adaptive Group Normalization (AdaGN) \\
Base Channels & 128 \\
Channel Mult. & (1, 1, 2, 2, 4, 4) for $256\times256$ inputs \\
Block Channels & (128, 128, 256, 256, 512, 512) \\
Residual Blocks & 2 per resolution level \\
Self-Attention & Enabled at $32\times32$, $16\times16$, $8\times8$ resolutions \\
Attention Blocks & \texttt{AttnDownBlock2D}, \texttt{AttnUpBlock2D} \\
\bottomrule
\end{tabular}
}%
}
\end{table}

\subsection{Source Image Conditioning via Channel Concatenation}
To condition on the source image $\bm{x}$, we adopt a channel-concatenation strategy. The input configuration is summarized in \cref{tab:suppl_conditioning}.

\begin{table}[h]
\centering
\caption{{\color{gray}Source image conditioning configuration.}}
\label{tab:suppl_conditioning}
{\color{gray}
\resizebox{\columnwidth}{!}{%
\begin{tabular}{ll}
\toprule
Component & Configuration \\
\midrule
Strategy & Channel-wise concatenation $[\bm{z}_t, \bm{x}]$ \\
Input Channels & $2C$ ($C$ is target channel count) \\
Output Channels & $C$ (raw denoised prediction) \\
\bottomrule
\end{tabular}
}%
}
\end{table}

\subsection{Bridge Scalings and Schedules}
Following the Diffusion Bridge framework, the raw model output $F_\theta$ is post-processed through bridge scalings $(c_{\text{skip}}, c_{\text{out}}, c_{\text{in}})$ derived from the VP log-SNR schedule. The denoised prediction is defined as:
\begin{equation}
    D_\theta(\bm{z}_t, t, \bm{x}) = c_{\text{out}}(t) \cdot F_\theta(c_{\text{in}}(t) \cdot \bm{z}_t,\, t,\, \bm{x}) + c_{\text{skip}}(t) \cdot \bm{z}_t\,,
    \label{eq:suppl_denoiser}
\end{equation}
The scheduling and numerical parameters are summarized in \cref{tab:suppl_bridge_params}.

\begin{table}[h]
\centering
\caption{{\color{gray}Bridge scheduling and numerical parameters.}}
\label{tab:suppl_bridge_params}
{\color{gray}
\resizebox{\columnwidth}{!}{%
\begin{tabular}{ll}
\toprule
Parameter & Value \\
\midrule
VP schedule $\beta_d$ & $2.0$ \\
VP schedule $\beta_{\min}$ & $0.1$ \\
Numerical stability $\varepsilon$ & $10^{-44}$ \\
Timestep rescaling & $\hat{t} = 1000 \cdot 0.25 \cdot \log(\sigma + \varepsilon)$ \\
\bottomrule
\end{tabular}
}%
}
\end{table}

\subsection{Training Hyperparameters}
The primary training objective is a weighted mean squared error between the denoised prediction $D_\theta(\bm{z}_t, t, \bm{x})$ and the clean target $\bm{y}$:
\begin{equation}
    \mathcal{L}_{\text{denoise}} = \mathbb{E}_{(\bm{x},\bm{y}),\, t,\, \epsilon} \left[ w(t) \cdot \| D_\theta(\bm{z}_t, t, \bm{x}) - \bm{y} \|_2^2 \right],
    \label{eq:suppl_denoising_loss}
\end{equation}
where $w(t) = A(t) / c_{\text{out}}(t)^2$ is the Karras bridge weighting. The diffusion time $t$ is sampled from a Karras distribution with parameters in \cref{tab:suppl_train_params}.

\begin{table}[h]
\centering
\caption{{\color{gray}Training time-sampling parameters.}}
\label{tab:suppl_train_params}
{\color{gray}
\resizebox{\columnwidth}{!}{%
\begin{tabular}{ll}
\toprule
Parameter & Value \\
\midrule
Distribution & Karras inverse-$\rho$ \\
$\rho$ factor & $7$ \\
Sampling range & $[\sigma_{\min}, \sigma_{\max}]$ \\
Concentration & Low noise levels (fine details) \\
\bottomrule
\end{tabular}
}%
}
\end{table}


\section{Additional Implementation Details}
\label{sec:suppl_impl}

{\color{gray}
This section provides a summary of the implementation hyperparameters and training environment used for EarthBridge.

\begin{table}[h]
\centering
\caption{{\color{gray}Implementation hyperparameters and training environment.}}
\label{tab:suppl_impl_params}
{\color{gray}
\resizebox{\columnwidth}{!}{%
\begin{tabular}{ll}
\toprule
Category & Details \\
\midrule
Optimizer & Prodigy (adaptive, nominal $lr=1.0$) \\
LR Schedule & Constant with 500 warmup steps \\
Gradient Clipping & Max norm $1.0$ \\
Mixed Precision & bfloat16 (Hugging Face Accelerate) \\
EMA Decay & $0.9999$ (with warmup) \\
Model Format & Hugging Face Diffusers (safetensors) \\
Multi-GPU & Distributed via Accelerate (7200s timeout) \\
Gradient Accum. & 1 (default) \\
\bottomrule
\end{tabular}
}%
}
\end{table}
}


\section{Dataset Processing Details}
\label{sec:suppl_dataset}

{\color{gray}
The MAVIC-T dataset is a composite of three geospatial data sources with standardized preprocessing and chipping. Details of the dataset processing are summarized in \cref{tab:suppl_dataset}.

\begin{table}[h]
\centering
\caption{{\color{gray}MAVIC-T dataset processing and augmentation details.}}
\label{tab:suppl_dataset}
{\color{gray}
\resizebox{\columnwidth}{!}{%
\begin{tabular}{ll}
\toprule
Component & Details \\
\midrule
Data Sources & UNICORN, USGS EROS HRO, UMBRA SAR \\
Preprocessing & Georectification, alignment, GeoTIFF storage \\
Chip Size & 200m$\times$200m spatially aligned stacks \\
Geographic Scope & UC Davis, Manhattan, Bingham, Centerfield \\
Augmentation & Random crop/resize (for 1024px tasks) \\
Quality Filter & Skip manifest (\texttt{bad\_samples.txt}) \\
\bottomrule
\end{tabular}
}%
}
\end{table}
}


\section{Per-Task Configuration Summary}
\label{sec:suppl_config}

\begin{table}[h]
\centering
\caption{{\color{gray}Per-task EarthBridge configuration summary.}}
\label{tab:task_config_suppl}
{\color{gray}
\resizebox{\columnwidth}{!}{%
\begin{tabular}{lcccc}
\toprule
Parameter & SAR$\rightarrow$EO & SAR$\rightarrow$RGB & SAR$\rightarrow$IR & RGB$\rightarrow$IR \\
\midrule
Source channels     & 1   & 1   & 1   & 3   \\
Target channels     & 1   & 3   & 1   & 1   \\
Model channels      & 1   & 3   & 1   & 3   \\
Resolution          & 256 & 1024 & 1024 & 1024 \\
Batch size          & 32  & 8   & 8   & 8   \\
Crop augmentation   & No  & Yes & Yes & Yes \\
Inference steps     & 40  & 40  & 40  & 40  \\
\bottomrule
\end{tabular}
}%
}
\end{table}

{\color{gray}
\cref{tab:task_config_suppl} summarizes the per-task configuration for our EarthBridge method. All tasks share the same UNet backbone architecture (128 base channels, channel multipliers auto-detected from resolution), Prodigy optimizer (lr=1.0), EMA decay (0.9999), VP noise schedule ($\sigma_{\min}=0.002$, $\sigma_{\max}=1.0$, $\beta_d=2.0$, $\beta_{\min}=0.1$), and training duration (100 epochs). The primary differences are the operating channel count, spatial resolution, and batch size.
}



\section{Additional Ablation Results}
\label{sec:suppl_ablation}

{\color{gray}
We provide additional ablation results to further characterize the behavior of the EarthBridge framework.
}

\paragraph{Inference Speed vs. Quality}
{\color{gray}
\cref{tab:inference_speed_suppl} reports EarthBridge inference time per image alongside the corresponding Task Score for the SAR$\rightarrow$EO task. As DBIM utilizes deterministic sampling, the generation quality remains remarkably stable even with extremely low NFE, making it suitable for time-sensitive remote sensing applications.
}

\begin{table}[h]
\centering
\caption{{\color{gray}EarthBridge inference speed and quality trade-off (SAR$\rightarrow$EO).}}
\label{tab:inference_speed_suppl}
{\color{gray}
\resizebox{\columnwidth}{!}{%
\begin{tabular}{lccc}
\toprule
Method & Steps (NFE) & Time (s) & Score$\downarrow$ \\
\midrule
EarthBridge & 10 & -- & -- \\
EarthBridge & 20 & -- & -- \\
EarthBridge & 40 & -- & -- \\
EarthBridge & 100 & -- & -- \\
\bottomrule
\end{tabular}
}
}
\end{table}

\paragraph{Multi-Stage Training Strategy}
{\color{gray}
Our two-stage training strategy is essential for handling the $1024\times1024$ tasks within memory constraints while capturing both local textures and global spatial dependencies. The strategy is summarized in \cref{tab:suppl_training_strategy}.

\begin{table}[h]
\centering
\caption{{\color{gray}Two-stage training strategy for high-resolution tasks.}}
\label{tab:suppl_training_strategy}
{\color{gray}
\resizebox{\columnwidth}{!}{%
\begin{tabular}{lll}
\toprule
Stage & Purpose & Details \\
\midrule
1: Pre-training & Local features & Random crop training on $1024^2$ data \\
2: Fine-tuning & Global structure & Full-resolution training on $1024^2$ images \\
\bottomrule
\end{tabular}
}%
}
\end{table}
}

\section{Additional Qualitative Results}
\label{sec:suppl_qualitative}

\begin{figure*}[t]
    \centering
    \fbox{\rule{0pt}{3in} \rule{0.95\textwidth}{0pt}}
    \caption{{\color{gray}Additional qualitative results for the SAR$\rightarrow$EO task. Left to right: source SAR image, EarthBridge output, ground truth. \TODO{Add visual comparison grid.}}}
    \label{fig:suppl_sar2eo}
\end{figure*}

\begin{figure*}[t]
    \centering
    \fbox{\rule{0pt}{3in} \rule{0.95\textwidth}{0pt}}
    \caption{{\color{gray}Additional qualitative results for the SAR$\rightarrow$RGB task. \TODO{Add visual comparison grid.}}}
    \label{fig:suppl_sar2rgb}
\end{figure*}

\begin{figure*}[t]
    \centering
    \fbox{\rule{0pt}{3in} \rule{0.95\textwidth}{0pt}}
    \caption{{\color{gray}Additional qualitative results for the RGB$\rightarrow$IR task. \TODO{Add visual comparison grid.}}}
    \label{fig:suppl_rgb2ir}
\end{figure*}

{\color{gray}
\cref{fig:suppl_sar2eo,fig:suppl_sar2rgb,fig:suppl_rgb2ir} provide additional qualitative results for three of the four translation tasks. \TODO{Add qualitative analysis.} These visualizations show a wider variety of scenes and input conditions, including challenging cases with complex urban structures, varying illumination, and dense vegetation.
\hypersetup{pageanchor=true}
}
